\section{Experimental protocol}

\subsection{Datasets and frozen evaluation splits}
We evaluate BioSNAP, Human and BindingDB under four scenarios: random,
cold-drug, cold-protein and cold-pair splits. Each scenario uses three fixed
random seeds and five outer folds, yielding 15 held-out evaluations per
dataset--scenario--method combination. Within each outer training partition,
approximately 10\% of eligible training positives are reserved for model
selection. For cold-drug and cold-protein splits, validation and test entities
are disjoint from training entities. For cold-pair splits, train, validation
and test positives are sampled from disjoint diagonal drug--protein blocks;
off-block pairs are not evaluated.

The three local \texttt{full.csv} files were verified as exact Git-blob matches
to the pinned DrugBAN dataset snapshot (commit
\texttt{9923f8c99959e00263103ff9ac61ba0eaccc8e02}). Its dataset documentation
identifies BindingDB, BioSNAP as distributed by MolTrans, and Human as
distributed by TransformerCPI \cite{bai2023drugban,liu2007bindingdb,huang2021moltrans,chen2020transformercpi}.
The generated dataset manifest records this source snapshot, byte-level
checksums and the observed input schema for every run.

For every partition, up to 10 unique unannotated pairs were sampled per
positive without replacement. Test and validation negatives were allocated
before training negatives to make held-out samples invariant to later
training-set changes. In unusually dense blocks with fewer available unique
negatives, all available unique negatives were used and the realized partition
size was retained in the fold-level result record.

\subsection{Realized benchmark coverage}
The completed release contains 180 frozen split definitions and 1,800
formal result records. This corresponds to 3 datasets, 4 split scenarios,
10 model variants and 15 seed--fold evaluations for every
dataset--scenario--method combination. All expected records passed the
release audit, and no dataset--scenario--method group had missing folds.
The realized test-set sizes were $30,426 \pm 0$ pairs for BioSNAP Random,
$5,793 \pm 6$ pairs for Human Random and $45,483 \pm 5$ pairs for
BindingDB Random. Under cold-pair evaluation, the corresponding test-set
sizes were $6,173 \pm 864$, $1,168 \pm 411$ and $9,107 \pm 1,732$ pairs,
respectively. The larger variation in the cold-start settings reflects the
number of eligible entities and unannotated drug--protein pairs available
within each frozen partition.

Across the complete benchmark, the mean number of test positives per fold
was $2,766 \pm 0$ for BioSNAP Random, $527 \pm 1$ for Human Random and
$4,135 \pm 0$ for BindingDB Random. In the cold-pair setting, the
corresponding values were $561 \pm 79$, $106 \pm 37$ and $828 \pm 157$.
These realized counts, together with the train and validation counts for
every dataset and split, are provided in Table~\ref{tab:fold-sizes} and in
the accompanying fold-level source-data CSV.

\subsection{Comparator families}
Table~\ref{tab:comparators} distinguishes feature-only, graph-based,
bilinear-attention and ablation controls. All rows use exactly the same frozen
positive partitions and partition-specific negative examples. ``Threshold GCN''
is a comparator implemented with the specified 0.3 within-type similarity
threshold; it is not presented as a full reproduction of an external method.
The panel is intended for controlled matched-partition comparisons rather than
an exhaustive reproduction of every published DTI method; conclusions are
therefore restricted to this benchmark.

\begin{table*}[t]
\centering
\scriptsize
\begin{tabular}{p{0.16\textwidth}p{0.22\textwidth}p{0.54\textwidth}}
\toprule
Family & Model & Information available to the model \\
\midrule
Feature-only & Logistic regression & Concatenated Morgan and ESM-2 pair features. \\
Feature-only & MLP & Concatenated Morgan and ESM-2 pair features. \\
Graph-based & GCN & Fold-local positive DTI graph with top-10 within-type similarity neighbours. \\
Graph-based & Threshold GCN & Fold-local positive DTI graph with within-type similarity edges above 0.3. \\
Bilinear-attention & DrugBAN & Canonical molecular graph and protein-sequence inputs, trained on the same frozen pair partitions. \\
Ablation & ArnoldiGCL w/o S1/S2; +S1; +S2; shuffled S1/S2 & The same fold-local ArnoldiGCL structural encoder, with the indicated side-information treatment. \\
Full model & ArnoldiGCL-S1/S2 & Structural pair features, raw Morgan and ESM-2 features, and aligned seven-dimensional S1/S2 information. \\
\bottomrule
\end{tabular}
\caption{Prespecified comparator families. All learning, graph construction
and label-derived side information are restricted to the outer training fold.}
\label{tab:comparators}
\end{table*}

\subsection{Evaluation and statistical reporting}
We report area under the receiver operating characteristic curve (AUC), area
under the precision--recall curve (AUPR), F1 score and accuracy. Checkpoints
were selected by validation AUC; F1 thresholds were selected by maximal
validation F1 and applied once to the test partition. The primary tables will
report mean $\pm$ sample standard deviation across 15 fixed seed--fold
evaluations. Pairwise comparisons will be paired by identical seed--fold
identity and summarized with effect sizes and two-sided Wilcoxon signed-rank
statistics. Because folds sharing a random seed are not independent biological
or experimental replicates, these $p$-values are descriptive summaries of
paired fold differences rather than confirmatory population-level inference.
For the complete family of reported descriptive comparisons, two-sided
$p$-values will be adjusted using the Benjamini--Hochberg procedure; both raw
and adjusted values will be provided as source data.

\subsection{Ablation and robustness analyses}
We compare the full S1/S2 model with a structural-only version, S1-only and
S2-only variants, and a partition-wise shuffled S1/S2 control. The shuffled
control distinguishes the presence of extra numerical inputs from alignment
between side information and the evaluated drug--protein pair.

\IfFileExists{generated/protocol_autogenerated.tex}{%
  \subsection{Fold-level dataset sizes}
Table~\ref{tab:fold-sizes} reports realized sample counts across the frozen folds. Counts may vary across cold-start folds because entities and eligible unannotated pairs are partition-specific.
\begin{table*}[t]
\centering
\scriptsize
\begin{tabular}{llrrrr}
\toprule
Dataset & Split & Train pairs & Validation pairs & Test pairs & Test positives \\
\midrule
BioSNAP & Random & 109538 $\pm$ 0 & 12166 $\pm$ 0 & 30426 $\pm$ 0 & 2766 $\pm$ 0 \\
BioSNAP & Cold drug & 109724 $\pm$ 1569 & 11980 $\pm$ 750 & 30426 $\pm$ 1336 & 2766 $\pm$ 121 \\
BioSNAP & Cold protein & 109120 $\pm$ 4026 & 12584 $\pm$ 3126 & 30426 $\pm$ 3839 & 2766 $\pm$ 349 \\
BioSNAP & Cold pair & 78327 $\pm$ 3469 & 1021 $\pm$ 315 & 6173 $\pm$ 864 & 561 $\pm$ 79 \\
\midrule
Human & Random & 20849 $\pm$ 6 & 2321 $\pm$ 0 & 5793 $\pm$ 6 & 527 $\pm$ 1 \\
Human & Cold drug & 21054 $\pm$ 989 & 2116 $\pm$ 912 & 5793 $\pm$ 1194 & 527 $\pm$ 109 \\
Human & Cold protein & 20789 $\pm$ 1099 & 2382 $\pm$ 843 & 5793 $\pm$ 1139 & 527 $\pm$ 104 \\
Human & Cold pair & 15101 $\pm$ 1453 & 150 $\pm$ 112 & 1168 $\pm$ 411 & 106 $\pm$ 37 \\
\midrule
BindingDB & Random & 163737 $\pm$ 5 & 18194 $\pm$ 0 & 45483 $\pm$ 5 & 4135 $\pm$ 0 \\
BindingDB & Cold drug & 163701 $\pm$ 700 & 18230 $\pm$ 466 & 45483 $\pm$ 685 & 4135 $\pm$ 62 \\
BindingDB & Cold protein & 163327 $\pm$ 8143 & 18605 $\pm$ 4576 & 45483 $\pm$ 8652 & 4135 $\pm$ 787 \\
BindingDB & Cold pair & 118036 $\pm$ 8799 & 1481 $\pm$ 538 & 9107 $\pm$ 1732 & 828 $\pm$ 157 \\
\bottomrule
\end{tabular}
\caption{Realized fold sizes (mean $\pm$ sample standard deviation across 15 fixed seed--fold evaluations).}
\label{tab:fold-sizes}
\end{table*}
%
}{%
  \todo{Insert final dataset statistics from the completed artifacts.}%
}

\subsection{Compute environment and reproducibility}
The principal benchmark was executed on NVIDIA GeForce RTX 4080 SUPER GPUs. Every
result artifact records the dataset, scenario, seed, outer fold, split
signature, validation criterion, test metrics and realized partition sizes.
DrugBAN artifacts additionally record the fixed 100-epoch schedule, input
preprocessing cache and TensorFloat-32 setting. The release-time environment
manifest records the software versions, GPU/driver information and hashes of
the key benchmark scripts. Aggregate wall-clock time is not reported because
per-fold duration is not a standardized field in every result artifact.
